% paper6_results.tex -- Results of Paper 6, input by papers/paper6_trials.tex.
% Every number is quoted from results/values.tex with \res{...}.
\section{Results}

The combined analysis is illustrated with maize grain yield from the nitrogen
rate trial of Paper~2, analysed across the three sites with
Equation~\eqref{eq:combined} (Table~\ref{tab:p6combined}).

\input{\subfix{../tables/p6_combined}}

Sites differed strongly ($F=\res{p6-F-site}$, \res{p6-P-site}), as expected
from their contrasting soils and climates, and nitrogen had a large effect
($F=\res{p6-F-nitrogen}$, \res{p6-P-nitrogen}). The site by nitrogen
interaction was also significant ($F=\res{p6-F-sitexnitrogen}$,
\res{p6-P-sitexnitrogen}): the size of the response to nitrogen differed
between sites, so the response should be reported site by site, as it is in
Paper~2, rather than as a single mean. Blocks within sites accounted for a
small but significant part of the variation ($F=\res{p6-F-blocksite}$,
\res{p6-P-blocksite}), which confirms that blocking removed real differences in
soil fertility from the error.

With \res{p6-df-error} residual degrees of freedom, $t_{0.025}=\res{p6-tcrit}$,
and the least significant difference between two site-by-rate means
(Equation~\eqref{eq:lsd}) is \res{p6-lsd}~\si{\tonne\per\hectare}. Two means
that differ by more than this amount differ at the 5 per cent level.
